\documentclass[12pt,a4paper]{article}

%% ============================================================
%% PACKAGES
%% ============================================================
\usepackage[utf8]{inputenc}
\usepackage{amsmath,amssymb,amsthm,mathtools}
\usepackage{geometry}
\usepackage{hyperref}
\usepackage{cleveref}
\usepackage{booktabs}
\usepackage{graphicx}
\usepackage{natbib}
\usepackage{tikz}
\usetikzlibrary{arrows.meta,positioning,calc}
\usepackage{algorithm}
\usepackage{algpseudocode}
\usepackage{enumitem}

\geometry{margin=1in}

\newtheorem{definition}{Definition}
\newtheorem{theorem}{Theorem}
\newtheorem{proposition}{Proposition}
\newtheorem{lemma}{Lemma}
\newtheorem{corollary}{Corollary}
\newtheorem{remark}{Remark}
\newtheorem{example}{Example}

\title{Mutual Liquidity Lock: A Bilateral Commitment Device \\
  with Graduated Enforcement on Blockchain}

\author{
  [Author Names] \\
  \textit{[Affiliations]} \\
  \texttt{[emails]}
}

\date{Working Paper --- March 2026}

\begin{document}
\maketitle

%% ============================================================
%% ABSTRACT
%% ============================================================

\begin{abstract}
We introduce the \textit{Mutual Liquidity Lock} (MLL), a trustless bilateral commitment protocol implemented as a smart contract. MLL enables two parties to make credible long-term commitments by locking assets in a shared pool with three novel enforcement mechanisms: (1) \textit{continuous commitment signals} via periodic deposits, (2) a \textit{bleeding penalty} that transfers value from a defaulting party to the compliant counterparty with a burned fraction, and (3) a \textit{graduated destruction mechanism} (``Russian Roulette'') that probabilistically freezes the pool with escalating stakes, implementing Schelling's brinkmanship in code. We provide a formal game-theoretic analysis showing that cooperation constitutes a Subgame Perfect Equilibrium for sufficiently patient players, with a critical discount factor $\delta^* = d/(d + r \cdot s)$ that \textit{decreases} as the pool matures---making the commitment self-reinforcing. We prove that the bleeding punishment is inherently renegotiation-resistant and that the Russian Roulette mechanism serves as both a deterrent and a bargaining catalyst. We analyze grief attacks, cold-start fragility, and behavioral deviations from rational play, providing design mitigations for each. A reference implementation in Solidity is provided. MLL represents a new financial primitive---bilateral liquidity lock---with applications to partnerships, service contracts, and any bilateral relationship requiring credible commitment without trusted third parties.

\medskip
\noindent\textbf{Keywords:} commitment devices, mechanism design, smart contracts, bilateral agreements, brinkmanship, game theory

\medskip
\noindent\textbf{JEL Classification:} C72, C73, D82, D86, G10
\end{abstract}

\newpage
\tableofcontents
\newpage

%% ============================================================
%% 1. INTRODUCTION
%% ============================================================

\section{Introduction}
\label{sec:intro}

The problem of credible commitment is fundamental to economics and social interaction. When two parties wish to engage in a long-term bilateral relationship---whether a business partnership, a service contract, or a personal commitment---each faces the risk that the other will defect. Traditional solutions include legal contracts (enforced by courts), third-party escrow (enforced by a trusted intermediary), and reputation systems (enforced by repeated interaction). Each has significant limitations: legal enforcement is slow, expensive, and jurisdiction-dependent; third-party escrow introduces counterparty risk; and reputation systems require transparency and repeated interaction in observable markets.

Blockchain technology and smart contracts offer a fundamentally different approach: \textit{algorithmic enforcement}. A commitment encoded in a smart contract executes automatically, without courts, intermediaries, or reputation. However, most existing smart contract mechanisms are designed for discrete transactions (e.g., atomic swaps, escrow) rather than ongoing bilateral relationships.

We propose the \textit{Mutual Liquidity Lock} (MLL) protocol, a mechanism that transforms abstract bilateral commitments into financially enforced, algorithmically executed agreements. MLL introduces three novel mechanisms that, in combination, create a robust commitment device:

\begin{enumerate}
  \item \textbf{Continuous Commitment Signals}: Both parties make periodic deposits into a shared pool, creating an observable ``heartbeat'' of commitment (\Cref{sec:continuous}).
  \item \textbf{Bleeding Penalty}: When one party ceases deposits, their share of the pool automatically drains---70\% to the counterparty, 30\% burned---converting passive defection into active financial loss (\Cref{sec:bleeding}).
  \item \textbf{Graduated Destruction (Russian Roulette)}: Either party can invoke a probabilistic mechanism that may freeze the pool for a configurable duration, implementing Schelling's brinkmanship \citep{schelling1960} with escalating stakes (\Cref{sec:roulette}).
\end{enumerate}

\subsection{Contributions}

Our contributions are threefold:

\textbf{Conceptual.} We identify bilateral liquidity lock as a new financial primitive---distinct from deposits, loans, insurance, or options---that enables credible commitment in trustless bilateral relationships. We connect this to the game-theoretic literature on commitment devices \citep{bryan2010}, hostage exchange \citep{williamson1983}, and brinkmanship \citep{schelling1960,powell1988}.

\textbf{Theoretical.} We provide a formal game-theoretic model of MLL as a discrete-time infinite-horizon repeated game. We prove that:
\begin{itemize}
  \item Mutual cooperation is a Subgame Perfect Equilibrium for sufficiently patient players, with a critical discount factor that decreases as the pool matures (\Cref{thm:cooperation}).
  \item The bleeding penalty is inherently renegotiation-resistant from the punisher's perspective (\Cref{prop:reneg}).
  \item The Russian Roulette mechanism serves as both a deterrent and a bargaining catalyst, driving disputing parties toward negotiated exit (\Cref{prop:rr-bargain}).
  \item In mature pools, depositing becomes a dominant strategy regardless of the opponent's action (\Cref{thm:dominant}).
  \item A share gate prevents grief attacks where low-stake parties weaponize graduated destruction (\Cref{prop:share-gate}).
\end{itemize}

\textbf{Practical.} We provide a reference implementation in Solidity, deployable on Ethereum L2 networks, with analysis of gas costs, security properties, and design mitigations for cold-start fragility and behavioral deviations.

\subsection{Related Work}
\label{sec:related}

\textbf{Commitment devices.} The theoretical foundations trace to \citet{schelling1960}, who showed that voluntarily constraining one's options can strengthen one's bargaining position. \citet{bryan2010} survey the modern literature, distinguishing hard commitments (binding constraints) from soft commitments (costly-to-break constraints). MLL combines both: the bleeding penalty is a soft commitment; the fund freeze is a hard commitment. Practical commitment platforms include stickK \citep{karlan2014} and Beeminder, but these are unilateral (self-commitment) rather than bilateral.

\textbf{Bilateral escrow and dual deposits.} \citet{asgaonkar2019} formalize dual-deposit escrow for one-shot digital good exchange, proving that honest cooperation is the unique Subgame Perfect Nash Equilibrium when deposits exceed the good's value. MLL extends this to the repeated-game setting with continuous deposits and graduated penalties.

\textbf{Hostage exchange.} \citet{williamson1983} introduces the concept of ``credible commitments using hostages to support exchange'' in transaction cost economics. MLL's mutual deposits are a direct implementation of Williamson's hostages, with the smart contract providing automatic enforcement that Williamson's framework assumed.

\textbf{Brinkmanship.} \citet{schelling1960,schelling1966} introduced the concept of ``the threat that leaves something to chance.'' \citet{powell1988,powell1990} formalized brinkmanship in models of nuclear deterrence. MLL's Russian Roulette mechanism is a direct code-level implementation of Schelling's brinkmanship, with the novel feature that the escalation probabilities are deterministic and publicly verifiable.

\textbf{Smart contract mechanism design.} \citet{cong2019} analyze how smart contracts expand the contracting space. The growing literature on DeFi mechanism design (see \citet{harvey2021}) provides the technological substrate for MLL's implementation.

\textbf{Contract theory.} \citet{hart1988} show that incomplete contracts with renegotiation constrain the feasible mechanism space. \citet{maskin1999} argue that optimal contracting is possible even with unforeseen contingencies if parties can commit not to renegotiate. MLL's immutable smart contract provides exactly this commitment.


%% ============================================================
%% 2. MODEL
%% ============================================================

\section{Model}
\label{sec:model}

\begin{definition}[MLL Game]
The \textbf{Mutual Liquidity Lock (MLL) Game} is a discrete-time, infinite-horizon, two-player game $\Gamma = \langle \{A, B\}, \mathcal{A}, \mathcal{S}, u, \delta \rangle$ defined as follows.
\end{definition}

\subsection{Players and Time}

Two players $i \in \{A, B\}$. Periods $t = 0, 1, 2, \ldots$ correspond to deposit intervals (e.g., 30 days). Both players share a common discount factor $\delta \in (0,1)$.

\subsection{State Space}

The game state at time $t$ is a tuple:
\[
\omega_t = (s_t^A, s_t^B, \beta_t, \kappa_t, \gamma_t, \phi_t)
\]
where:
\begin{itemize}
    \item $s_t^i \geq 0$: player $i$'s \textit{share} (claim on the pool), with total pool $P_t = s_t^A + s_t^B$
    \item $\beta_t \in \{0, A, B\}$: bleeding target ($0$ = no bleeding active)
    \item $\kappa_t \in \{0, 1, \ldots, 6\}$: Russian Roulette invocation counter
    \item $\gamma_t \in \{0, 1, \ldots, \Theta\}$: cooldown timer ($\Theta$ periods after RR invocation)
    \item $\phi_t \in \{\texttt{active}, \texttt{frozen}, \texttt{exited}\}$: game phase
\end{itemize}

\textbf{Initial state:} $\omega_0 = (s_0, s_0, 0, 0, 0, \texttt{active})$ where $s_0 > 0$ is each player's initial deposit.

\subsection{Action Space}

When $\phi_t = \texttt{active}$ and $\gamma_t = 0$ (no cooldown), each player $i$ simultaneously chooses:
\[
a_t^i \in \mathcal{A} = \{D, S, R, E, U\}
\]
\begin{itemize}
    \item $D$ (Deposit): contribute the agreed amount $d > 0$
    \item $S$ (Skip): do not contribute
    \item $R$ (Roulette): invoke the Russian Roulette mechanism (requires deposit of $d$ and $\geq 35\%$ pool share)
    \item $E$ (Exit): propose peaceful bilateral dissolution
    \item $U$ (Unilateral Exit): initiate unilateral withdrawal with penalty
\end{itemize}

During cooldown ($\gamma_t > 0$): $\mathcal{A} = \{D, S, E, U\}$ (no RR).

During warmup (first $W$ periods): $\mathcal{A} = \{D, S, E\}$ (no RR or unilateral exit).

When $\phi_t \in \{\texttt{frozen}, \texttt{exited}\}$: no actions available (terminal).

\subsection{Transition Rules}
\label{sec:transitions}

\subsubsection{Deposit and Bleeding}

Given action profile $(a_t^A, a_t^B)$:

\textbf{Case 1: Both deposit} ($a_t^A = a_t^B = D$).
\[
s_{t+1}^i = s_t^i + d, \quad \beta_{t+1} = 0
\]

\textbf{Case 2: Player $i$ deposits, player $j$ skips} ($a_t^i = D, a_t^j = S$), and $j$ has exceeded the grace period $G$:

Let $r_t$ denote the effective bleeding rate at time $t$ (see \Cref{sec:decaying-rate}). The bleed amount is $\Delta_t = r_t \cdot s_t^j$. Of this:
\begin{align}
s_{t+1}^i &= s_t^i + d + \alpha \cdot \Delta_t \label{eq:bleed-to-counterparty} \\
s_{t+1}^j &= s_t^j - \Delta_t \label{eq:bleed-from-defaulter}
\end{align}
where $\alpha = 0.70$ is the counterparty fraction and $(1-\alpha) \cdot \Delta_t$ is burned (permanently removed from both parties' claims). The burn creates deadweight loss that reduces spite incentives while maintaining the punisher's incentive to enforce (\Cref{sec:bleed-split}).

\textbf{Case 3: Both skip} ($a_t^A = a_t^B = S$).
No bleeding transfer (symmetric default):
$s_{t+1}^i = s_t^i, \quad \beta_{t+1} = 0$.

\begin{remark}
Symmetric default produces no directional transfer. This prevents a party from claiming victimhood while also not depositing.
\end{remark}

\subsubsection{Peaceful Exit}

If both players choose $E$: $\phi_{t+1} = \texttt{exited}$. Each player receives their current share $s_t^i$.

\subsubsection{Unilateral Exit}

If player $i$ chooses $U$, a countdown of $T_U$ periods begins. During the countdown, bleeding continues. After the countdown:
\begin{align}
\text{Player } i \text{ receives:} &\quad (1 - p_U) \cdot s_{T_U}^i \label{eq:unilateral-exit} \\
\text{Player } j \text{ receives:} &\quad s_{T_U}^j + p_U \cdot s_{T_U}^i
\end{align}
where $p_U$ is the exit penalty (default 15\%). The game transitions to $\phi = \texttt{exited}$.

\subsubsection{Russian Roulette Mechanism}
\label{sec:rr-def}

If player $i$ chooses $R$ (requires $\gamma_t = 0$, $t > W$, deposit of $d$, and share $\geq 35\%$ of pool):

\begin{enumerate}
    \item Counter increments: $\kappa_{t+1} = \kappa_t + 1$
    \item \textbf{Trigger check}: With probability $p(\kappa_{t+1}) = \frac{1}{7 - \kappa_{t+1}}$, funds freeze for duration $T_F$:
    $\phi_{t+1} = \texttt{frozen}$.
    With probability $1 - p(\kappa_{t+1})$, the game continues with cooldown:
    $\gamma_{t+1} = \Theta$.
    \item Cooldown decrements each period: $\gamma_{t+1} = \max(\gamma_t - 1, 0)$
\end{enumerate}

The probability sequence:
\begin{center}
\begin{tabular}{c|cccccc}
$k$ (invocation) & 1 & 2 & 3 & 4 & 5 & 6 \\
\hline
$p(k)$ & $\frac{1}{6}$ & $\frac{1}{5}$ & $\frac{1}{4}$ & $\frac{1}{3}$ & $\frac{1}{2}$ & $1$ \\
Cumulative & $\frac{1}{6}$ & $\frac{1}{3}$ & $\frac{1}{2}$ & $\frac{2}{3}$ & $\frac{5}{6}$ & $1$
\end{tabular}
\end{center}

\begin{remark}
This is equivalent to drawing without replacement from $\{1,2,3,4,5,6\}$. Each invocation is equally likely to be the fatal one ex ante ($\Pr[\text{trigger at exactly invocation } k] = 1/6$ for all $k$), while the conditional probability monotonically increases.
\end{remark}

\subsection{Terminal Payoffs}

\textbf{Frozen.} When $\phi = \texttt{frozen}$, funds are locked for $T_F$ periods (default 10 years). After $T_F$ expires, each player receives their share at freeze time. Since $\delta^{T_F} \approx 0$ for realistic discount factors and $T_F = 3650$ days, the effective loss to both parties is approximately total.

\textbf{Peaceful exit.} Each player receives $s_{\text{exit}}^i$.

\textbf{Unilateral exit.} Initiator receives $(1-p_U) \cdot s^i$; counterparty receives $s^j + p_U \cdot s^i$.


%% ============================================================
%% 3. CONTINUOUS COMMITMENT SIGNALS
%% ============================================================

\section{Continuous Commitment Signals}
\label{sec:continuous}

\subsection{Deposits as Costly Signals}

In Spence's signaling framework \citep{spence1973}, a signal is credible if and only if its cost is higher for the type that would benefit from mimicry. In MLL, the periodic deposit $d$ functions as a costly signal of continued commitment: a party intending to defect gains nothing from depositing (the deposit increases their exposure to loss), while a committed party benefits from the growing pool.

Formally, the signaling cost differential between committed type $\theta_C$ and defecting type $\theta_D$ is:
\[
c(\theta_C) - c(\theta_D) = d - d = 0 \text{ (myopically)}
\]
but in the multi-period setting:
\[
c(\theta_C, T) - c(\theta_D, T) = -\sum_{t=1}^{T} \delta^t \cdot r \cdot s_t^j < 0
\]

The committed type faces \textit{lower} cost because they avoid bleeding. This ensures the signal is separating: only parties intending to continue would deposit.

\subsection{Heartbeat Semantics}

Each deposit serves as a ``heartbeat'' confirming liveness and intent:

\begin{enumerate}
\item \textbf{Liveness proof}: The deposit transaction proves the party still controls their private key and is monitoring the contract.
\item \textbf{Intent signal}: By increasing their locked exposure, the depositor credibly signals continued commitment.
\item \textbf{State synchronization}: Like a TCP keepalive, the deposit cadence provides a clock for the bilateral relationship, enabling automated detection of counterparty failure.
\end{enumerate}

The grace period $G$ (default: 1 missed interval) provides tolerance for operational failures (e.g., gas spikes, temporary key unavailability) without triggering punishment.

\subsection{Dead Man's Switch}

If a party fails to interact with the contract for $T_{\text{dead}} = 90$ days, the counterparty gains unilateral claim to the entire pool. This addresses the key scenario, ensuring that key loss by one party does not permanently lock the counterparty's funds.


%% ============================================================
%% 4. BLEEDING PENALTY
%% ============================================================

\section{The Bleeding Penalty}
\label{sec:bleeding}

The bleeding penalty is MLL's primary enforcement mechanism. When one party ceases deposits while the counterparty continues, the defaulter's share automatically drains at a calibrated rate. This section formalizes the bleeding dynamics, analyzes the 70/30 split design, and proves renegotiation-resistance.

\subsection{Bleeding Dynamics}

Let player $j$ default (play $S$) while player $i$ continues depositing (play $D$). After the grace period $G$, bleeding activates. Player $j$'s share evolves as:
\[
s_{t+\tau}^j = s_t^j \cdot (1 - r)^\tau
\]
where $r$ is the per-period bleeding rate and $\tau$ counts periods since bleeding activation.

The total bleed per period is $\Delta_\tau = r \cdot s_{t+\tau-1}^j$. This is split:
\begin{itemize}
\item $\alpha \cdot \Delta_\tau$ added to counterparty $i$'s share
\item $(1-\alpha) \cdot \Delta_\tau$ burned (removed from both claims)
\end{itemize}

\subsection{The 70/30 Split Design}
\label{sec:bleed-split}

The split parameter $\alpha$ faces a fundamental tension:

\textbf{$\alpha = 1$ (100\% to counterparty):} The punisher actively benefits from punishment, making the punishment self-enforcing. However, this creates a \textit{spite incentive}: the punisher may refuse reasonable exit offers because continued bleeding is more profitable. The mechanism degenerates into extraction.

\textbf{$\alpha = 0$ (100\% burned):} The punisher gains nothing from enforcing punishment. Since the punisher must continue depositing to sustain the bleeding (otherwise symmetric default → no bleeding), punishment becomes costly with zero reward. The punisher becomes indifferent between punishing and exiting, and rational punishment collapses to at most one period. Deterrence is destroyed.

\begin{proposition}[Lower Bound on $\alpha$]
\label{prop:alpha-bound}
For the bleeding punishment to be individually rational for the punisher, $\alpha$ must satisfy:
\[
\alpha > \frac{d}{r \cdot s_t^j}
\]
i.e., the punisher's share of the bleed transfer must exceed the per-period deposit cost.
\end{proposition}

\begin{proof}
The punisher's per-period payoff from continuing punishment is:
\[
\pi_{\text{punish}} = \alpha \cdot r \cdot s_t^j - d
\]
For punishment to be individually rational, $\pi_{\text{punish}} \geq 0$, giving $\alpha \geq d / (r \cdot s_t^j)$. For mature pools where $s_t^j \gg d$, this constraint is easily satisfied even for small $\alpha$.
\end{proof}

Our default $\alpha = 0.70$ provides a large margin above the lower bound while burning 30\% to limit spite. The burned fraction creates deadweight loss that makes both parties worse off during punishment, aligning incentives toward peaceful resolution.

\subsection{Decaying Bleeding Rate (Cold Start Mitigation)}
\label{sec:decaying-rate}

At pool inception, the accumulated stake $s_0 = d$ (a single deposit). The critical discount factor (\Cref{thm:cooperation}):
\[
\delta^*(t=0) = \frac{d}{d + r \cdot d} = \frac{1}{1 + r}
\]

For typical $r \approx 0.15$ per month, $\delta^* \approx 0.87$---uncomfortably close to 1 for moderately impatient agents. The cooperation equilibrium barely holds; a small external shock could trigger defection.

To address this cold-start fragility, we introduce a decaying bleeding rate:
\[
r_t = \begin{cases}
r_\infty \cdot \left(2 - \frac{t}{T_{\text{decay}}}\right) & \text{if } t \leq T_{\text{decay}} \\
r_\infty & \text{if } t > T_{\text{decay}}
\end{cases}
\]
where $r_\infty$ is the steady-state rate and $T_{\text{decay}}$ is the decay period (default: 6 deposit intervals). At inception, $r_0 = 2 \cdot r_\infty$, decaying linearly to $r_\infty$.

This yields:
\[
\delta^*(t=0) = \frac{d}{d + 2r_\infty \cdot d} = \frac{1}{1 + 2r_\infty} \approx 0.77
\]
substantially lowering the cooperation threshold during the vulnerable early phase.

\subsection{Renegotiation-Resistance}
\label{sec:reneg-bleeding}

\begin{proposition}[Bleeding is Renegotiation-Resistant]
\label{prop:reneg}
The bleeding punishment is self-enforcing: in every punishment period, the punisher strictly prefers to continue punishment over any renegotiation that reduces the punishment.
\end{proposition}

\begin{proof}
When $B$ is bleeding and $A$ is the punisher, $A$'s continuation value from bleeding is:
\[
V_A^{\text{bleed}} = \sum_{\tau=1}^{\infty} \delta^\tau \left[\alpha \cdot r \cdot s_t^B (1-r)^{\tau-1} - d\right]
= \frac{\delta(\alpha \cdot r \cdot s_t^B - d)}{1 - \delta(1-r)}
\]
This is positive when $\alpha \cdot r \cdot s_t^B > d$, which holds in any non-trivial pool.

For $A$ to accept renegotiation to peaceful exit, $A$ must receive at least:
\[
s_t^A + V_A^{\text{bleed}} > s_t^A
\]
But $B$ can only offer $B$'s own remaining share $s_t^B$. When $V_A^{\text{bleed}} > s_t^B$ (i.e., the present value of future bleeding exceeds $B$'s remaining balance), no exit offer from $B$ can satisfy $A$. This occurs when:
\[
\frac{\delta \cdot \alpha \cdot r \cdot s_t^B}{1 - \delta(1-r)} > s_t^B \implies \frac{\delta \cdot \alpha \cdot r}{1 - \delta(1-r)} > 1
\]
which simplifies to $\alpha > \frac{1-\delta}{\delta \cdot r}$---easily satisfied for patient players with reasonable $r$.
\end{proof}

\begin{remark}[Trust Barrier]
Formally, MLL fails BPW renegotiation-proofness \citep{hart1988}: both parties could agree off-chain to exit and re-enter a fresh agreement. However, renegotiation requires \textit{bilateral trust}---precisely the trust that failed when punishment was triggered. When the mechanism is most needed (low trust), renegotiation is least feasible. This ``trust barrier'' provides de facto renegotiation-resistance.
\end{remark}


%% ============================================================
%% 5. GRADUATED DESTRUCTION
%% ============================================================

\section{Graduated Destruction: The Russian Roulette Mechanism}
\label{sec:roulette}

\subsection{Brinkmanship in Code}

\citet{schelling1960} defined brinkmanship as ``the deliberate creation of a recognizable risk of war, a risk that one does not completely control.'' The Russian Roulette (RR) mechanism implements this precisely: invoking RR creates a verifiable risk of pool freeze that neither party can retract once initiated.

Key properties of the implementation:
\begin{enumerate}
\item \textbf{Verifiable randomness}: Production deployment uses Chainlink VRF, providing publicly auditable randomness. Neither party can influence the outcome.
\item \textbf{Escalating stakes}: Each invocation increases the trigger probability ($1/6 \to 1/5 \to \cdots \to 1$), creating accelerating urgency.
\item \textbf{Guaranteed termination}: At most 6 invocations before certain trigger, bounding the brinkmanship horizon.
\item \textbf{Cooldown}: 10-day pause between invocations prevents impulsive escalation and provides a window for negotiation.
\end{enumerate}

\subsection{Formal Analysis}

\begin{proposition}[RR as Bargaining Catalyst]
\label{prop:rr-bargain}
In the Russian Roulette subgame with counter $\kappa = k$, if both players are risk-averse (concave utility), there exists a peaceful exit split $(\alpha, 1-\alpha)$ that both players prefer to continuing the RR process.
\end{proposition}

\begin{proof}
Let $u(\cdot)$ be a concave utility function. Player $i$'s expected utility from continuing RR:
\[
EU_i^{\text{RR}}(k) = p(k+1) \cdot u(u_{\text{frozen}}^i) + (1-p(k+1)) \cdot EU_i^{\text{RR}}(k+1)
\]
where $u_{\text{frozen}}^i \approx 0$ (effective total loss at realistic $T_F$). By concavity:
\[
u(\alpha \cdot P) > \alpha \cdot u(P) + (1-\alpha) \cdot u(0) \geq EU_i^{\text{RR}}
\]
for some $\alpha \in (0,1)$, since the certainty equivalent of the RR lottery is strictly less than $P$ for any risk-averse agent. As $p(k)$ increases, the certainty equivalent decreases, widening the bargaining range.
\end{proof}

The key insight is that RR does not need to trigger to be effective: the escalating probability drives convergence to a negotiated exit.

\begin{proposition}[RR Deterrence Threshold]
\label{prop:rr-deterrence}
Player $i$ refrains from invoking RR if and only if:
\[
V_{\text{status quo}}^i \geq (1 - p(k+1)) \cdot V_{\text{continue}}^i + p(k+1) \cdot u_{\text{frozen}}^i
\]
A player invokes RR when the status quo (e.g., continued bleeding) has eroded their position enough that the gamble becomes preferable.
\end{proposition}

\subsection{Grief Attack Mitigation}
\label{sec:grief}

A critical vulnerability: a party with negligible pool share could weaponize RR to freeze a much larger counterparty's funds. The cost to the attacker (loss of their small share) is far less than the damage (freezing the victim's large share).

\begin{proposition}[Share Gate Prevents Grief Attacks]
\label{prop:share-gate}
Requiring $s_t^i / P_t \geq \eta$ (with $\eta = 0.35$) to invoke RR ensures that any invocation imposes cost proportional to the damage.
\end{proposition}

\begin{proof}
Under the share gate, the invoker's expected loss from a trigger is:
\[
L_{\text{invoker}} = p(k) \cdot s_t^i \geq p(k) \cdot \eta \cdot P_t = \eta \cdot p(k) \cdot P_t
\]
The counterparty's expected loss is:
\[
L_{\text{counter}} = p(k) \cdot s_t^j \leq p(k) \cdot (1-\eta) \cdot P_t
\]
The cost ratio $L_{\text{invoker}} / L_{\text{counter}} \geq \eta / (1-\eta)$. For $\eta = 0.35$, this is $\geq 0.538$: the invoker bears at least 54\% of the cost, making grief attacks unprofitable.

Moreover, the share gate is self-reinforcing: if the griefer's share falls below $\eta$ due to bleeding, they lose RR access precisely when the attack would be most asymmetric.
\end{proof}

Additional mitigations:
\begin{itemize}
\item \textbf{Warmup period}: First $W = 3$ deposit intervals, no RR access. Prevents day-1 grief attacks.
\item \textbf{Deposit-to-invoke}: Must deposit $d$ to invoke RR, creating a per-invocation cost floor.
\end{itemize}


%% ============================================================
%% 6. MECHANISM PROPERTIES
%% ============================================================

\section{Mechanism Properties}
\label{sec:properties}

\subsection{Subgame Perfect Equilibrium}

\begin{theorem}[Cooperation Sustainability]
\label{thm:cooperation}
The cooperative strategy profile $\sigma^*$ (both players deposit in every period) constitutes a Subgame Perfect Equilibrium if and only if:
\[
\delta \geq \delta^* \equiv \frac{d}{d + r \cdot s_t^i}
\]
where $s_t^i$ is player $i$'s current share at the point of potential deviation.
\end{theorem}

\begin{proof}
Consider player $B$'s incentive to deviate from $D$ to $S$ at time $t$, given that $A$ follows $\sigma^*$.

\textbf{One-period gain from deviation}: $+d$ (saved deposit).

\textbf{Cost of deviation}: Bleeding activates. $B$'s share decays at rate $r$ per period. The present value of the loss, assuming $B$ deviates permanently:
\[
\text{PV}_{\text{bleed}} = \sum_{\tau=1}^{\infty} \delta^\tau \cdot r \cdot s_t^B \cdot (1-r)^{\tau-1}
= \frac{\delta \cdot r \cdot s_t^B}{1 - \delta(1-r)}
\]

For deviation to be unprofitable: $d \leq \text{PV}_{\text{bleed}}$. When $r \ll 1$ and $\delta$ is close to 1:
\[
\text{PV}_{\text{bleed}} \approx \frac{\delta \cdot r \cdot s_t^B}{1 - \delta + \delta r} \approx \frac{r \cdot s_t^B}{(1-\delta)/\delta + r}
\]

Setting $d = \text{PV}_{\text{bleed}}$ and solving for $\delta$:
\[
\delta \geq \delta^* = \frac{d}{d + r \cdot s_t^B}
\]
\end{proof}

\begin{corollary}[Pool Maturity Effect]
\label{cor:maturity}
$\delta^*$ is strictly decreasing in $P_t$. The MLL becomes more stable over time: the longer parties have cooperated, the harder it is to justify deviation.
\end{corollary}

\begin{example}
With $d = 100$, $s = 10{,}000$ (after extended cooperation), $r = 0.15$:
$\delta^* = 100 / (100 + 1500) = 0.0625$.
Even agents discounting the future at 93.75\% would find cooperation rational.
\end{example}

\subsection{Dominant Strategy in Mature Pools}

\begin{theorem}[Deposit Dominance]
\label{thm:dominant}
When $r \cdot s_t^j > d / \delta$, playing $D$ is a dominant strategy regardless of the opponent's action.
\end{theorem}

\begin{proof}
If opponent plays $D$: $D$ yields share growth $d$; $S$ saves $d$ but risks bleeding upon opponent retaliation.

If opponent plays $S$: $D$ yields the bleeding transfer $\alpha \cdot r \cdot s_t^j$ minus deposit cost $d$. Since $\alpha \cdot r \cdot s_t^j > d$ when $r \cdot s_t^j > d/\alpha$, the net gain is positive. Playing $S$ yields zero transfer (symmetric default).

In both cases, the payoff from $D$ exceeds that from $S$.
\end{proof}

\subsection{Individual Rationality}

\begin{theorem}[Ex-Ante Individual Rationality]
\label{thm:ir}
Participation in MLL is individually rational when the commitment value $v$ that cooperation produces (per period) satisfies:
\[
v > d \cdot \frac{1-\delta}{\delta}
\]
i.e., the per-period value of the bilateral relationship exceeds the opportunity cost of the locked deposit.
\end{theorem}

This is the ``it's worth it if the relationship is worth it'' condition: MLL does not create value by itself but amplifies the value of bilateral commitments by making them credible.

\subsection{Unilateral Exit Preserves Properties}

\begin{proposition}[Unilateral Exit as Legal Escape Valve]
\label{prop:unilateral}
The unilateral exit mechanism (countdown $T_U$ + penalty $p_U$) preserves the cooperation equilibrium while eliminating legal vulnerability:
\begin{enumerate}[label=(\roman*)]
\item The total unilateral exit cost is approximately $p_U + r \cdot T_U \approx 22.5\%$ of share (for defaults: $p_U = 15\%$, $r \cdot T_U \approx 7.5\%$). This is high enough to preserve deterrence.
\item The existence of unilateral exit prevents the contract from being challenged as unconscionable (no exit possible).
\item The cooperation equilibrium $\delta^*$ is unaffected: unilateral exit is off the equilibrium path.
\end{enumerate}
\end{proposition}


%% ============================================================
%% 7. PARAMETER DESIGN
%% ============================================================

\section{Parameter Design and Calibration}
\label{sec:parameters}

\subsection{Bleeding Rate $r$}

The bleeding rate must satisfy three constraints:

\begin{enumerate}
    \item \textbf{Deterrence}: $r > d(1-\delta) / (s_t^i \cdot \delta)$ (deviation is unprofitable).
    \item \textbf{Recovery}: $(1-r)^G > \alpha_{\min}$ where $G$ is the grace period and $\alpha_{\min} \approx 0.80$ is minimum share retention during grace.
    \item \textbf{RR calibration}: The number of periods until RR invocation becomes rational should be meaningful but not excessive: $T_{\text{RR}} = \ln(\theta_{\text{RR}}) / \ln(1-r)$.
\end{enumerate}

\begin{proposition}[Optimal Bleeding Rate]
For baseline parameters ($d = 0.1$ ETH monthly, $s_0 = 1$ ETH, $\delta \approx 0.99$), the optimal range is:
\[
r^* \in [0.003, 0.008] \text{ per day} \quad \Leftrightarrow \quad [30, 80] \text{ bps/day}
\]
Our default $r = 50$ bps/day ($\approx 14\%$ per month) falls within this range.
\end{proposition}

\subsection{Russian Roulette Parameters}

\begin{proposition}[Optimality of the $1/(7-k)$ Sequence]
The sequence $p(k) = 1/(7-k)$ is equivalent to sampling without replacement from $\{1,\ldots,6\}$. It has two desirable properties:
\begin{enumerate}
    \item \textbf{Uniform marginal}: $\Pr[\text{trigger at exactly invocation } k] = 1/6$ for all $k$.
    \item \textbf{Monotone hazard rate}: $p(k)$ is strictly increasing, ensuring escalation.
\end{enumerate}
\end{proposition}

Alternative: i.i.d. draws with fixed $p$ (geometric distribution). This has constant hazard rate and does not guarantee termination. The without-replacement design is strictly preferable for brinkmanship.

\subsection{Parameter Summary}

\begin{table}[h]
\centering
\begin{tabular}{llll}
\toprule
\textbf{Parameter} & \textbf{Symbol} & \textbf{Default} & \textbf{Rationale} \\
\midrule
Deposit interval & $\Delta t$ & 30 days & Monthly cadence \\
Bleeding rate & $r$ & 50 bps/day & $\sim$14\%/month deterrence \\
Bleed split & $\alpha$ & 70/30 & Punisher IC + spite reduction \\
Grace period & $G$ & 1 interval & Operational tolerance \\
Warmup & $W$ & 3 intervals & Anti-grief \\
Decay period & $T_{\text{decay}}$ & 6 intervals & Cold start mitigation \\
RR cooldown & $\Theta$ & 10 days & Negotiation window \\
RR invocations & $K_{\max}$ & 6 & Bounded brinkmanship \\
Share gate & $\eta$ & 35\% & Grief prevention \\
Freeze duration & $T_F$ & 10 years & $\approx$ permanent at realistic $\delta$ \\
Exit penalty & $p_U$ & 15\% & Deterrence + legal escape \\
Exit countdown & $T_U$ & 30 days & Response window \\
Dead man's switch & $T_{\text{dead}}$ & 90 days & Key loss mitigation \\
\bottomrule
\end{tabular}
\caption{MLL parameter defaults with rationale.}
\label{tab:params}
\end{table}


%% ============================================================
%% 8. IMPLEMENTATION
%% ============================================================

\section{Implementation}
\label{sec:implementation}

\subsection{Architecture}

The reference implementation consists of two Solidity contracts:

\begin{enumerate}
\item \texttt{RussianRoulette.sol}: Abstract contract implementing the RR mechanism with cooldown management, probability calculation, and pseudo-random trigger (Chainlink VRF for production).
\item \texttt{MutualLiquidityLock.sol}: Main contract inheriting \texttt{RussianRoulette}, implementing the full MLL protocol including deposit management, bleeding calculation, 70/30 split, share gate, warmup, decaying rate, unilateral exit, and configurable freeze.
\end{enumerate}

\subsection{Bleeding Computation}

Bleeding is computed incrementally on each interaction (deposit, withdrawal, or explicit call). For a defaulting party with share $s$ and last bleed timestamp $t_0$, the bleed amount at time $t_1$ is:
\[
\Delta = r_{\text{eff}} \cdot s \cdot \frac{t_1 - t_0}{1 \text{ day}}
\]
where $r_{\text{eff}}$ is the effective rate from \Cref{sec:decaying-rate}. This linear approximation avoids the gas cost of computing $(1-r)^\tau$ on-chain while remaining accurate for typical interaction frequencies.

\subsection{Security Considerations}

\begin{itemize}
\item \textbf{Randomness}: MVP uses \texttt{block.prevrandao}; production requires Chainlink VRF v2.5 to prevent miner manipulation.
\item \textbf{Reentrancy}: ETH transfers use checks-effects-interactions pattern. No external contract calls during state mutation.
\item \textbf{Overflow}: Solidity 0.8.24+ with built-in overflow checks. Share arithmetic uses 256-bit integers.
\item \textbf{Front-running}: RR invocation is front-runnable (miner could skip blocks to influence \texttt{prevrandao}). Chainlink VRF eliminates this.
\item \textbf{Griefing}: Share gate (\Cref{sec:grief}) prevents low-cost grief attacks on the RR mechanism.
\end{itemize}

\subsection{Gas Costs}

Estimated gas costs (Ethereum mainnet, Solidity 0.8.24 with IR optimizer):
\begin{center}
\begin{tabular}{lr}
\toprule
\textbf{Operation} & \textbf{Gas} \\
\midrule
Activation (per party) & $\sim$50K \\
Deposit & $\sim$35K \\
Invoke RR & $\sim$65K \\
Peaceful exit & $\sim$55K \\
Unilateral exit (initiate) & $\sim$30K \\
Unilateral exit (execute) & $\sim$60K \\
\bottomrule
\end{tabular}
\end{center}

At typical L2 costs ($<$\$0.01 per transaction), MLL is operationally negligible. On Ethereum mainnet, monthly deposits at 50 gwei would cost $\sim$\$2--5, which is small relative to typical deposit amounts.


%% ============================================================
%% 9. APPLICATIONS
%% ============================================================

\section{Applications}
\label{sec:applications}

\subsection{Business Partnerships}

Two parties entering a joint venture can use MLL to credibly commit to continued collaboration. The periodic deposit models ongoing capital contribution, the bleeding penalty deters silent abandonment, and the RR mechanism provides an escalation path if the partnership deteriorates. The unilateral exit with penalty allows either party to leave (at a cost), analogous to a buyout clause.

\subsection{Long-Term Service Contracts}

A client and service provider can lock mutual deposits. The client commits to continued payment (deposits); the provider commits to continued service. Default by either party triggers bleeding, providing automatic compensation without litigation. This is particularly valuable for cross-border agreements where legal enforcement is impractical.

\subsection{DAO-to-DAO Agreements}

Decentralized organizations that wish to formalize bilateral cooperation (e.g., liquidity sharing, protocol integration) can use MLL as an on-chain commitment device. The trustless nature eliminates the need for legal wrappers around DAO-to-DAO agreements.

\subsection{Commitment for Personal Relationships}

While unconventional, MLL could formalize financial commitments in personal relationships (e.g., cohabitation agreements, collaborative projects). The mechanism provides a middle ground between informal trust and formal legal contracts. Adverse selection concerns apply: parties who seek MLL may be in already-strained relationships. However, as with couples therapy, selection on need does not preclude benefit.


%% ============================================================
%% 10. DISCUSSION
%% ============================================================

\section{Discussion}
\label{sec:discussion}

\subsection{Behavioral Considerations}

The formal model assumes rational, risk-neutral agents with exponential discounting. Real agents deviate in several ways that affect MLL's dynamics:

\textbf{Loss aversion.} Prospect theory predicts that losses loom larger than equivalent gains \citep{kahneman1979}. In MLL, this amplifies the deterrent effect of bleeding (losses feel $\sim$2x as large) but also predicts that players will trigger RR \textit{earlier} than the rational model suggests, as the bleeding loss triggers disproportionate emotional response. The 10-day cooldown is more important than the probability sequence for preventing impulsive escalation.

\textbf{Hyperbolic discounting (dual role).} Present-biased agents have effectively higher $\delta^*$, making cooperation harder to sustain. However, MLL is simultaneously a \textit{solution to} present bias: the protocol itself is a commitment device against one's future impatient self, in the tradition of Ulysses binding himself to the mast. MLL is weakened by and a cure for the same cognitive bias---a duality worth formalizing.

\textbf{Sunk cost fallacy.} Agents tend to over-weight sunk costs. In MLL, accumulated deposits create a sunk cost that rationally should be ignored when evaluating future actions---but behaviorally anchors agents to continued participation. This ``bug'' reinforces the rational cooperation equilibrium: agents stay committed partly because they have ``already invested so much.'' A behavioral bias that acts as a feature.

\subsection{Limitations}

\begin{itemize}
\item \textbf{Capital inefficiency}: Locked funds earn no yield in the v1 implementation. Yield integration (v2) would address this but adds composability risk.
\item \textbf{Correlation risk}: Both parties' commitment is denominated in the same asset (ETH). If ETH price crashes, both parties' outside options change simultaneously, potentially destabilizing the equilibrium.
\item \textbf{Key management}: Loss of a private key leaves funds locked until the dead man's switch triggers (90 days). Multisig or social recovery could mitigate this.
\item \textbf{Regulatory uncertainty}: Whether MLL constitutes a financial instrument under securities law is jurisdiction-dependent and untested.
\end{itemize}

\subsection{Comparison to Legal Contracts}

\begin{center}
\begin{tabular}{lcc}
\toprule
\textbf{Property} & \textbf{MLL} & \textbf{Legal Contract} \\
\midrule
Enforcement speed & Immediate & Months/years \\
Enforcement cost & Gas fee & Legal fees \\
Jurisdiction & Global & Local \\
Flexibility & Rigid (code) & Flexible (interpretation) \\
Renegotiation & Hard (trust barrier) & Easy (amendment) \\
Completeness & Complete & Incomplete \\
Exit option & Expensive ($\sim$22.5\%) & Varies \\
\bottomrule
\end{tabular}
\end{center}

MLL and legal contracts are complements, not substitutes. MLL excels in cross-border, low-trust, or legally sparse environments; legal contracts excel where flexibility and contextual interpretation are needed.


%% ============================================================
%% 11. CONCLUSION
%% ============================================================

\section{Conclusion}
\label{sec:conclusion}

We have introduced the Mutual Liquidity Lock, a bilateral commitment protocol that implements credible commitment through three interlocking mechanisms: continuous commitment signals, bleeding penalties with a 70/30 split, and graduated destruction via Russian Roulette. The formal analysis shows that cooperation is a self-reinforcing equilibrium that strengthens as the pool matures, with a critical discount factor $\delta^* = d/(d + r \cdot s)$ that approaches zero for mature pools.

The key theoretical insight is the mechanism's self-reinforcing nature: each period of cooperation increases the accumulated stake, which lowers $\delta^*$, which makes cooperation easier to sustain, which leads to further cooperation. Once the pool passes an initial vulnerability threshold (mitigated by the decaying bleeding rate), it becomes increasingly robust.

The design space explored here---bilateral liquidity locks with graduated enforcement---is largely unexplored in both the game theory and DeFi literatures. We identify several directions for future work:

\begin{enumerate}
\item \textbf{Yield integration}: Enabling the locked pool to earn DeFi yield would address capital inefficiency but introduces composability risk and changes the equilibrium analysis.
\item \textbf{Multi-party extension}: $N > 2$ players with coalition dynamics, connecting to team production theory \citep{holmstrom1982}.
\item \textbf{Adaptive parameters}: On-chain parameter adjustment based on pool state, potentially using bonding curves for the bleeding rate.
\item \textbf{Empirical validation}: Deploying MLL on testnet/mainnet and analyzing real behavioral patterns against the theoretical predictions.
\item \textbf{Formal verification}: Mechanized proofs of the smart contract's correctness properties.
\end{enumerate}


%% ============================================================
%% REFERENCES
%% ============================================================

\bibliographystyle{plainnat}

\begin{thebibliography}{99}

\bibitem[Asgaonkar and Krishnamachari(2019)]{asgaonkar2019}
Asgaonkar, A. and Krishnamachari, B. (2019).
\newblock Solving the buyer and seller's dilemma: A dual-deposit escrow smart contract for provably cheat-proof delivery and payment for a digital good without a trusted mediator.
\newblock In \textit{IEEE International Conference on Blockchain and Cryptocurrency (ICBC)}.

\bibitem[Bryan et~al.(2010)]{bryan2010}
Bryan, G., Karlan, D., and Nelson, S. (2010).
\newblock Commitment devices.
\newblock \textit{Annual Review of Economics}, 2:671--698.

\bibitem[Cong and He(2019)]{cong2019}
Cong, L.~W. and He, Z. (2019).
\newblock Blockchain disruption and smart contracts.
\newblock \textit{Review of Financial Studies}, 32(5):1754--1797.

\bibitem[Grossman and Hart(1986)]{grossman1986}
Grossman, S.~J. and Hart, O.~D. (1986).
\newblock The costs and benefits of ownership: A theory of vertical and lateral integration.
\newblock \textit{Journal of Political Economy}, 94(4):691--719.

\bibitem[Hart and Moore(1988)]{hart1988}
Hart, O. and Moore, J. (1988).
\newblock Incomplete contracts and renegotiation.
\newblock \textit{Econometrica}, 56(4):755--785.

\bibitem[Harvey et~al.(2021)]{harvey2021}
Harvey, C.~R., Ramachandran, A., and Santoro, J. (2021).
\newblock \textit{DeFi and the Future of Finance}.
\newblock John Wiley \& Sons.

\bibitem[Holmstr{\"o}m(1982)]{holmstrom1982}
Holmstr{\"o}m, B. (1982).
\newblock Moral hazard in teams.
\newblock \textit{Bell Journal of Economics}, 13:324--340.

\bibitem[Kahneman and Tversky(1979)]{kahneman1979}
Kahneman, D. and Tversky, A. (1979).
\newblock Prospect theory: An analysis of decision under risk.
\newblock \textit{Econometrica}, 47(2):263--291.

\bibitem[Karlan and Appel(2014)]{karlan2014}
Karlan, D. and Appel, J. (2014).
\newblock Failing in the field: What we can learn when field research goes wrong.
\newblock \textit{Journal of Economic Perspectives}.

\bibitem[Maskin and Tirole(1999)]{maskin1999}
Maskin, E. and Tirole, J. (1999).
\newblock Unforeseen contingencies and incomplete contracts.
\newblock \textit{Review of Economic Studies}, 66(1):83--114.

\bibitem[Powell(1988)]{powell1988}
Powell, R. (1988).
\newblock Nuclear brinkmanship with two-sided incomplete information.
\newblock \textit{American Political Science Review}, 82(1):155--178.

\bibitem[Powell(1990)]{powell1990}
Powell, R. (1990).
\newblock \textit{Nuclear Deterrence Theory: The Search for Credibility}.
\newblock Cambridge University Press.

\bibitem[Schelling(1960)]{schelling1960}
Schelling, T.~C. (1960).
\newblock \textit{The Strategy of Conflict}.
\newblock Harvard University Press.

\bibitem[Schelling(1966)]{schelling1966}
Schelling, T.~C. (1966).
\newblock \textit{Arms and Influence}.
\newblock Yale University Press.

\bibitem[Spence(1973)]{spence1973}
Spence, M. (1973).
\newblock Job market signaling.
\newblock \textit{Quarterly Journal of Economics}, 87(3):355--374.

\bibitem[Williamson(1983)]{williamson1983}
Williamson, O.~E. (1983).
\newblock Credible commitments: Using hostages to support exchange.
\newblock \textit{American Economic Review}, 73(4):519--540.

\end{thebibliography}

\end{document}
